\chapter{Tests de résistance et scénarios}\label{ch:rc:stress-testing-and-scenarios}

En mars 2021, la chute des cours des positions les plus concentrées d'Archegos Capital
Management, un family office, déclenche des appels de marge qu'il ne peut pas honorer. En un an, il était passé
d'environ 1.5~milliard de dollars de capital et de 10 milliards d'exposition à plus de 36 milliards de
capital et 160 milliards d'exposition, en grande partie au moyen de total return swaps conclus avec plusieurs prime
brokers, dont chacun ne voyait que sa propre part. Son défaut provoque plus de 10 milliards de
dollars de pertes dans plusieurs grandes banques~; l'une d'elles, Credit Suisse, perd près de 5.5~milliards de dollars. Ces pertes proviennent d'un
mouvement qu'une mesure de risque calibrée sur le passé récent, appliquée à la tranche d'un seul courtier
du portefeuille, n'aurait pas jugé probable. Ce chapitre porte sur les outils qui posent l'autre
question~: non pas combien le portefeuille perd un mauvais jour, mais ce qui se passe dans un monde défavorable particulier,
quels mondes feraient tomber l'établissement, et ce que les superviseurs exigent des banques pour survivre.

\section{Ce que la valeur en risque ne voit pas}

\begin{definition}[Test de résistance]\label{def:rc:stress-testing-and-scenarios:stress}
Un \emph{test de résistance}\index{test de résistance} réévalue un portefeuille, ou l'ensemble d'un établissement, selon un
scénario sévère mais plausible de mouvements de \omterm{def:rc:market-risk-measures:factor}{facteurs de risque} défini à l'avance, et en rapporte la perte et ses
conséquences (fonds propres, liquidité, limites), sans attacher de probabilité au scénario.
\end{definition}

La \omterm{def:rc:market-risk-measures:var}{valeur en risque} (VaR) et l'ES résument une distribution estimée sur une fenêtre d'historique~; elles ne voient pas les mouvements que la
fenêtre ne contient pas, les changements de corrélation, la concentration et l'illiquidité d'une
position (qui allongent le délai de sortie), ni les pertes des contreparties qui détiennent les mêmes
positions. Un \omterm{def:rc:stress-testing-and-scenarios:stress}{test de résistance} remplace la distribution par un récit, raconté en chiffres.

\section{Scénarios historiques}

\begin{definition}[Scénario historique]\label{def:rc:stress-testing-and-scenarios:historical}
Un \emph{scénario historique}\index{scénario historique} applique aux positions d'aujourd'hui les mouvements que les
\omterm{def:rc:market-risk-measures:factor}{facteurs de risque} ont connus sur une période passée~: différences pour les taux et les spreads, variations relatives
pour les prix.
\end{definition}

\begin{example}[Trois épisodes sur un même portefeuille]\label{ex:rc:stress-testing-and-scenarios:historical}
Le portefeuille du chapitre~21 (dix ans longs sur le Trésor américain, deux ans courts, euros longs, yens courts, un
straddle EURUSD vendeur), réévalué intégralement sur trois périodes de données du Trésor américain et de la BCE~:
\begin{itemize}
\item Lehman, du 12 septembre au 10 octobre 2008~: rendement à deux ans de $-61$ points de base, à dix ans $+15$,
EURUSD $-3.52~\%$, USDJPY $-7.85~\%$~; une perte de 17.26~millions de dollars~;
\item mars 2020, du 6 au 20 mars~: $-12$ et $+18$ points de base, EURUSD $-5.71~\%$, USDJPY $+5.11~\%$~; une
perte de 20.13~millions de dollars~;
\item septembre 2022, du 21 au 28 septembre~: $+5$ et $+21$ points de base, EURUSD $-3.50~\%$, USDJPY
$+0.46~\%$~; une perte de 11.61~millions de dollars.
\end{itemize}
La VaR historique à dix jours et 99~\% issue du chapitre~21 s'élève à 6.63~millions de dollars~: chaque épisode fait perdre deux
à trois fois plus (\cref{fig:rc:stress-testing-and-scenarios:historical}). Le straddle vendeur,
inoffensif les jours ordinaires, domine les pertes dans les trois cas.
\end{example}

\begin{omfigure}
\begin{tikzpicture}
\begin{axis}[width=11.5cm,height=5.2cm,xbar,bar width=12pt,
  symbolic y coords={ten-day VaR,September 2022,March 2020,Lehman 2008},ytick=data,yticklabels={{VaR à dix jours},{septembre 2022},{mars 2020},{Lehman 2008}},
  xlabel={perte (millions de dollars)},
  tick label style={font=\small},label style={font=\small},
  xmin=0,xmax=24,grid=major,grid style={gray!25},
  nodes near coords,every node near coord/.append style={font=\footnotesize,/pgf/number format/fixed,/pgf/number format/precision=1},
  table/col sep=comma]
\addplot[fill=omThm!60,draw=omThm] table[y=scenario,x=loss]{figdata/rates-credit-risk/22-stress-testing-and-scenarios/historical.csv};
\end{axis}
\end{tikzpicture}

\omcaption{Pertes du portefeuille du chapitre~21 selon trois \omterm{def:rc:stress-testing-and-scenarios:historical}{scénarios historiques}, à comparer à sa VaR à dix jours
à 99~\%. Chaque épisode coûte deux à trois fois la VaR. Données~: taux au pair du Trésor américain, taux
de référence de la BCE~; tutoriel du chapitre.}\label{fig:rc:stress-testing-and-scenarios:historical}
\end{omfigure}

\begin{dated}{2026-09}{Octobre 1987 et septembre 2022}\label{dat:rc:stress-testing-and-scenarios:history}
Le 19 octobre 1987, le Dow Jones Industrial Average perd 508 points, soit 22.6~\%, en une seule séance, la plus forte
baisse d'un jour de son histoire. Fin septembre 2022, le rendement du gilt britannique à trente ans grimpe de
160 points de base en quelques jours, avec une amplitude intrajournalière de 127 points de base sur une même journée, les fonds
en gestion adossée au passif (LDI) ayant dû faire face à des appels de marge.
\end{dated}

\section{Scénarios hypothétiques et leur cohérence}

\begin{definition}[Scénario hypothétique]\label{def:rc:stress-testing-and-scenarios:hypothetical}
Un \emph{scénario hypothétique}\index{scénario hypothétique} est un ensemble de mouvements conçus plutôt qu'observés~:
un récit (une flambée des taux, une dévaluation, un défaut) traduit en chocs sur certains \omterm{def:rc:market-risk-measures:factor}{facteurs de risque}, le reste des
facteurs étant fixé de manière cohérente.
\end{definition}

\begin{definition}[Conditionnement de scénario]\label{def:rc:stress-testing-and-scenarios:conditioning}
Le \emph{conditionnement de scénario}\index{conditionnement de scénario} remplit les facteurs que le scénario ne précise
pas par leur espérance conditionnelle sachant ceux qui sont précisés. Sous un modèle gaussien conjoint de covariance $\Sigma$, les chocs $x_s$
sur les facteurs précisés donnent $\E[x_u\mid x_s] = \Sigma_{us}\Sigma_{ss}^{-1}x_s$ pour les
autres.
\end{definition}

\omcode{code/firm/stresstest/firm_stresstest.py}{32}{40}{Conditionnement~: les facteurs qu'un scénario laisse ouverts, à leur espérance conditionnelle.}\label{lst:rc:stress-testing-and-scenarios:condition}

\begin{example}[Un choc sur dix ans]\label{ex:rc:stress-testing-and-scenarios:hypo}
Choquer le rendement à dix ans à la hausse de 100 points de base. Seul, il coûte au portefeuille 14.80~millions de dollars. Complété
par la covariance à dix jours des mouvements quotidiens depuis 2020, le scénario fait aussi monter le rendement à deux ans de 80
points de base, baisser l'euro de 0.36~\% et monter l'USDJPY de 2.25~\%~; la perte est de 10.04~millions de dollars, car les
deux ans courts gagnent. Le scénario se situe à 5.3 écarts types du centre de la
distribution à dix jours.
\end{example}

\section{Tests de résistance inversés}

\begin{definition}[Test de résistance inversé]\label{def:rc:stress-testing-and-scenarios:reverse}
Un \emph{test de résistance inversé}\index{test de résistance inversé} part d'un résultat, une perte qui ferait franchir une limite,
les fonds propres ou l'établissement lui-même, et cherche les scénarios qui la produisent~; sa forme quantitative trouve
celui qui est le plus plausible, le mouvement de plus petite distance de Mahalanobis
$\sqrt{x^\top\Sigma^{-1}x}$ avec une perte d'au moins $L$.
\end{definition}

\begin{proposition}[Le scénario de perte le plus plausible d'un portefeuille linéaire]\label{prop:rc:stress-testing-and-scenarios:rev}
Pour un P\&L $\delta^\top x$ et une covariance $\Sigma$, le mouvement de plus petite distance de Mahalanobis avec une perte $L$ est
$x^* = -L\,\Sigma\delta/(\delta^\top\Sigma\delta)$, à la distance $L/\sqrt{\delta^\top\Sigma\delta}$~: la perte mesurée en écarts types
du P\&L du portefeuille.
\end{proposition}
\begin{proof}
Minimiser $\tfrac12x^\top\Sigma^{-1}x$ sous la contrainte $\delta^\top x = -L$~: la condition de stationnarité du lagrangien donne $\Sigma^{-1}x = \mu\delta$,
donc $x = \mu\Sigma\delta$, et la contrainte fixe $\mu$.
\end{proof}

Pour un portefeuille non linéaire, la méthode itère~: on prend la direction $\Sigma\nabla\mathrm{P\&L}$ au point courant,
on se déplace le long de cette direction jusqu'à ce que la perte visée soit atteinte, puis on recalcule le gradient en ce point.

\begin{example}[Franchir une limite de 10~millions d'USD]\label{ex:rc:stress-testing-and-scenarios:rev}
Avec la covariance à dix jours, le scénario le plus plausible du modèle linéaire pour une perte de 10~millions d'USD
fait monter le rendement à deux ans de 33 points de base et celui à dix ans de 58, fait baisser l'euro de 3.18~\% et fait monter
l'USDJPY de 1.78~\%, à 3.93 écarts types~; réévalué en totalité, ce scénario fait perdre 14.00
millions d'USD. La recherche en réévaluation complète en trouve un plus proche, à 2.47 écarts types~: l'euro
baisse de 3.45~\%, l'USDJPY monte de 2.15~\%, les rendements montent de 10 et 17 points de base
(voir la \cref{fig:rc:stress-testing-and-scenarios:reverse}). Le straddle vendu rend les fortes variations de devises,
dans un sens comme dans l'autre, le moyen le moins cher de perdre.
\end{example}

\begin{omfigure}
\begin{tikzpicture}
\begin{axis}[width=11.5cm,height=5.2cm,ybar,bar width=11pt,
  symbolic x coords={2y,10y,EURUSD,USDJPY},xtick=data,xticklabels={2A,10A,EURUSD,USDJPY},
  ylabel={mouvement (écarts types à dix jours)},
  tick label style={font=\small},label style={font=\small},
  ymin=-3.5,ymax=3.5,grid=major,grid style={gray!25},
  legend columns=2,legend style={font=\footnotesize,at={(0.5,-0.2)},anchor=north,draw=none,fill=none,/tikz/every even column/.append style={column sep=8pt}},
  table/col sep=comma]
\addplot[fill=omDef!60,draw=omDef] table[x=factor,y=linear]{figdata/rates-credit-risk/22-stress-testing-and-scenarios/reverse.csv};
\addplot[fill=omThm!60,draw=omThm] table[x=factor,y=full]{figdata/rates-credit-risk/22-stress-testing-and-scenarios/reverse.csv};
\legend{modèle linéaire (3.93 é.-t.),réévaluation complète (2.47 é.-t.)}
\end{axis}
\end{tikzpicture}

\omcaption{Les scénarios à dix jours les plus plausibles qui font perdre 10~millions d'USD au portefeuille, issus du modèle
linéaire et de la réévaluation complète, chaque mouvement de facteur étant exprimé dans ses propres écarts types à dix jours.
Le modèle linéaire répartit le mouvement entre taux et devises~; avec la convexité du straddle, une variation de devise
fait les dégâts à une distance bien plus petite. Données~: le tutoriel
du chapitre.}\label{fig:rc:stress-testing-and-scenarios:reverse}
\end{omfigure}

\section{Tests de résistance prudentiels}

\begin{definition}[Test de résistance prudentiel]\label{def:rc:stress-testing-and-scenarios:supervisory}
Un \emph{test de résistance prudentiel}\index{test de résistance prudentiel} (supervisory stress test) est un \omterm{def:rc:stress-testing-and-scenarios:stress}{test de résistance} dont les scénarios
sont fixés par le superviseur et appliqués à de nombreuses banques à la fois, afin de dimensionner les coussins de fonds propres et de comparer
la résilience des banques.
\end{definition}

Les scénarios de la Réserve fédérale américaine combinent des trajectoires macroéconomiques et financières avec, pour les banques
dotées de grandes activités de trading, un choc de marché global appliqué à leurs positions de trading et autres positions valorisées
à la juste valeur et, pour certaines, une composante de défaut de contrepartie. L'Autorité bancaire européenne organise
des tests à l'échelle de l'UE selon une méthodologie commune.

\begin{dated}{2026-09}{Scénarios prudentiels}\label{dat:rc:stress-testing-and-scenarios:fed}
Le scénario très défavorable définitif 2026 de la Réserve fédérale (publié en février 2026) est une récession
mondiale sévère déclenchée par une baisse de l'appétit pour le risque~: le taux de chômage augmente de 5.5 points
de pourcentage jusqu'à un pic de 10~\% au troisième trimestre 2027, les cours des actions baissent d'environ 58~\% jusqu'au
troisième trimestre 2026, les prix immobiliers résidentiels de 30~\% et ceux de l'immobilier commercial de 39~\%, et le VIX atteint un pic
à 72. Dans le \omterm{def:rc:stress-testing-and-scenarios:stress}{test de résistance} de l'Autorité bancaire européenne de 2025, à l'échelle de l'UE (64 banques, 75~\% des actifs bancaires de l'UE), le
scénario défavorable a coûté 370 points de base de fonds propres de catégorie~1 sous forme d'actions ordinaires, laissant un ratio global de 12~\%~; la
méthodologie provisoire de l'exercice 2027 a été publiée le 11 juin 2026.
\end{dated}

\section{Tutoriel~: une bibliothèque de scénarios}\label{tut:rc:stress-testing-and-scenarios:tutorial}

\begin{tutorial}
\textbf{Objectif.} Appliquer des \omterm{def:rc:stress-testing-and-scenarios:historical}{scénarios historiques} au portefeuille du chapitre~21, compléter un \omterm{def:rc:stress-testing-and-scenarios:hypothetical}{scénario hypothétique} par
conditionnement, et trouver le scénario le plus plausible derrière une perte donnée. \textbf{État final~:} les
chiffres des
\cref{ex:rc:stress-testing-and-scenarios:historical,ex:rc:stress-testing-and-scenarios:hypo,ex:rc:stress-testing-and-scenarios:rev}
et les graphiques.
\begin{tutsteps}
\item \textbf{Bibliothèque}~: \texttt{scenarios()} construit les trois épisodes avec \texttt{historical}.
\item \textbf{Exécution}~: \texttt{historical\_table()}~; comparer avec \texttt{var10()}.
\item \textbf{Conditionnement}~: \texttt{hypothetical(100)}.
\item \textbf{Inversion}~: \texttt{reverse(10e6)}~; \texttt{fig\_rc\_stress.py} écrit les graphiques.
\end{tutsteps}
\textbf{Pour aller plus loin.} Conditionner sur la covariance stressée de mars 2020 à la place~; mettre le
straddle à zéro et répéter le \omterm{def:rc:stress-testing-and-scenarios:reverse}{test de résistance inversé}.
\end{tutorial}

\section{Construction~: le moteur de test de résistance}\label{bld:rc:stress-testing-and-scenarios:stresstest}

\begin{build}
\bfield{Rôle} La bibliothèque de scénarios de la banque et les calculs de stress, exécutés chaque nuit par le moteur
de risque du chapitre~29 en même temps que la VaR.
\bfield{Interface} \texttt{Scenario(name, moves)}~; \texttt{historical(name, dates, levels,
start, end, log\_cols)}~; \texttt{condition(cov, shocked, values)}~; \texttt{mahalanobis}~;
\texttt{reverse\_linear(delta, cov, loss)}~; \texttt{reverse\_nonlinear(pnl, cov, loss, x0)}~;
\texttt{run(pnl, scenarios)}.
\bfield{Règles} Mouvements dans les unités du modèle (rendements en décimales, prix en log-rendements)~;
plausibilité par la distance de Mahalanobis sous la covariance retenue.
\bfield{Tests d'acceptation} \texttt{code/firm/stresstest/tests/}~: mouvements historiques~; le conditionnement
égale la régression~; le scénario inverse linéaire se situe sur le plan de perte et est plus proche que tout autre
point de ce plan~; la recherche non linéaire le retrouve pour un portefeuille linéaire.
\bfield{Pour aller plus loin} Covariances stressées et plausibilité à queues épaisses~; bibliothèques de scénarios par
régime~; horizons ajustés de la liquidité par position~; \omterm{def:rc:stress-testing-and-scenarios:reverse}{test de résistance inversé} à l'échelle de la firme, incluant le financement.
\end{build}

\begin{omsources}
\item US Securities and Exchange Commission, communiqué de presse 2022-70, 27 avril 2022.
\item Paul, Weiss, rapport adressé au Special Committee of the Board of Credit Suisse Group sur
Archegos Capital Management, 29 juillet 2021.
\item Federal Reserve History, «~Stock market crash of 1987~».
\item Board of Governors of the Federal Reserve System, \emph{2026 Stress Test Scenarios}, 2026.
\item Bank of England, lettre de Sir Jon Cunliffe adressée au Treasury Committee, 5 octobre 2022.
\end{omsources}

\section{Exercices}

\begin{exercise}[$\star$]\label{exo:rc:stress-testing-and-scenarios:1}
Pourquoi une perte de stress peut-elle être plusieurs fois supérieure à la VaR sans que le modèle de VaR soit faux~?
\end{exercise}

\begin{exercise}[$\star$]\label{exo:rc:stress-testing-and-scenarios:2}
Dans un modèle normal à deux facteurs, de volatilités 10 et 5 points de base et de corrélation 0.6, le
premier facteur subit un choc de 40 points de base. Quel est le mouvement conditionnel du second~?
\end{exercise}

\begin{exercise}[$\star$]\label{exo:rc:stress-testing-and-scenarios:3}
Un portefeuille linéaire présente un écart-type de P\&L à dix jours de 2.5~millions de dollars. À quelle
distance de Mahalanobis perd-il 10~millions de dollars~?
\end{exercise}

\begin{exercise}[$\star\star$]\label{exo:rc:stress-testing-and-scenarios:4}
Pourquoi le choc hypothétique de dix ans a-t-il coûté moins cher une fois complété par conditionnement~?
\end{exercise}

\begin{exercise}[$\star\star$]\label{exo:rc:stress-testing-and-scenarios:5}
Pourquoi le scénario linéaire du \omterm{def:rc:stress-testing-and-scenarios:reverse}{test de résistance inversé} n'est-il pas le plus plausible pour ce livre~?
\end{exercise}

\begin{exercise}[$\star\star$]\label{exo:rc:stress-testing-and-scenarios:6}
Que suppose un \omterm{def:rc:stress-testing-and-scenarios:historical}{scénario historique} sur les positions et les marchés d'aujourd'hui qui pourrait être faux~?
\end{exercise}

\begin{exercise}[$\star\star\star$]\label{exo:rc:stress-testing-and-scenarios:7}
\emph{Programmation.} Exécuter le \omterm{def:rc:stress-testing-and-scenarios:reverse}{test de résistance inversé} pour une perte de 10~millions de dollars sans le straddle et
comparer la distance.
\end{exercise}

\begin{exercise}[$\star\star\star$]\label{exo:rc:stress-testing-and-scenarios:8}
\emph{Trouver la faille.} «~Notre scénario de \omterm{def:rc:stress-testing-and-scenarios:reverse}{test de résistance inversé} se situe à 4 écarts-types, donc il
ne se produit qu'une fois sur 30\,000 périodes de dix jours et nous pouvons le négliger.~»
\end{exercise}

\section{Problème~: le family office vu du côté du courtier}

\begin{problem}[{Problème du week-end --- quand la marge s'épuise}]\label{pb:rc:stress-testing-and-scenarios:1}
Un prime broker fait face à un client au travers de total return swaps portant sur cinq actions concentrées, pour un notionnel
de 20~milliards de dollars pondéré à 30, 25, 20, 15 et 10~\%, couvert par la détention des actions. Le client verse
une marge de 7.5~\% du notionnel. Si le client fait défaut, le courtier a besoin de cinq jours pour vendre. Dans un
modèle illustratif à un facteur, chaque action présente un bêta de 1.2 par rapport à un marché dont la volatilité quotidienne est de 1.2~\%
et une volatilité idiosyncratique quotidienne de 3~\%.

\medskip\noindent\textbf{Partie I --- Le livre.}
\begin{enumerate}
\item Donner la marge en dollars.
\item Donner l'écart-type de la valeur des actions sur cinq jours.
\item Pourquoi le risque du courtier est-il le défaut du client et non les actions~?
\item Que dirait une VaR à un jour à 99~\% des swaps pendant que le client verse sa marge~?
\item Qu'est-ce que le courtier n'a pas vu et que détenaient les autres courtiers du client~?
\end{enumerate}

\medskip\noindent\textbf{Partie II --- Stress inversé.}
\begin{enumerate}[resume]
\item Donner les mouvements sur cinq jours les plus plausibles qui consomment la marge.
\item Donner leur distance de Mahalanobis et la probabilité normale d'un mouvement au moins aussi éloigné.
\item Pourquoi les positions les plus importantes chutent-elles le plus dans ce scénario~?
\item Quelle baisse uniforme consomme la marge~?
\item Qu'est-ce qui change avec une marge de 20~\%~?
\end{enumerate}

\medskip\noindent\textbf{Partie III --- Au-delà du modèle.}
\begin{enumerate}[resume]
\item Pourquoi cinq jours seraient-ils optimistes~?
\item Comment les ventes des autres courtiers modifient-elles le scénario~?
\item Pourquoi le taux de marge des swaps était-il lui-même un risque~?
\item Quelle information exigerait une marge dynamique~?
\item Quelle limite de concentration sur un seul nom faudrait-il fixer~?
\end{enumerate}

\medskip\noindent\textbf{Partie IV --- Jugement.}
\begin{enumerate}[resume]
\item Pourquoi s'agit-il d'un problème de \omterm{def:rc:stress-testing-and-scenarios:stress}{test de résistance} et non d'un problème de VaR~?
\item Comment le résultat du \omterm{def:rc:stress-testing-and-scenarios:reverse}{test de résistance inversé} doit-il entrer dans les décisions du courtier~?
\item Que faudrait-il demander au client avant d'accorder davantage de capacité~?
\item Énoncer le \emph{résultat nommé}~: les mouvements auxquels le livre épuise sa marge, et
leur distance.
\item En une phrase~: que dit un \omterm{def:rc:stress-testing-and-scenarios:reverse}{test de résistance inversé} à un comité des risques~?
\end{enumerate}
\end{problem}

\section{Questions d'entretien}

\begin{interviewq}[$\star$ \iqroles{risque}]\label{iq:rc:stress-testing-and-scenarios:1}
Quelle est la différence entre une VaR et un \omterm{def:rc:stress-testing-and-scenarios:stress}{test de résistance}~? Quand feriez-vous confiance à chacun~?
\end{interviewq}

\begin{interviewq}[$\star\star$ \iqroles{risque, chercheur}]\label{iq:rc:stress-testing-and-scenarios:2}
Comment rendez-vous cohérent en interne un \omterm{def:rc:stress-testing-and-scenarios:hypothetical}{scénario hypothétique}~?
\end{interviewq}

\begin{interviewq}[$\star\star$ \iqroles{chercheur}]\label{iq:rc:stress-testing-and-scenarios:3}
Dérivez le scénario le plus plausible pour une perte donnée d'un portefeuille linéaire.
\end{interviewq}

\begin{interviewq}[$\star\star$ \iqroles{banque, risque}]\label{iq:rc:stress-testing-and-scenarios:4}
Que teste le global market shock de la Réserve fédérale, et en quoi diffère-t-il de la VaR~?
\end{interviewq}

\begin{interviewq}[$\star\star\star$ \iqroles{risque, trader}]\label{iq:rc:stress-testing-and-scenarios:5}
Comment soumettriez-vous à un \omterm{def:rc:stress-testing-and-scenarios:stress}{test de résistance} un livre de prime brokerage~?
\end{interviewq}

\begin{interviewq}[$\star\star\star$ \iqroles{développeur}]\label{iq:rc:stress-testing-and-scenarios:6}
Concevez une bibliothèque de scénarios pour le passage nocturne des \omterm{def:rc:stress-testing-and-scenarios:stress}{tests de résistance} d'une banque. Que contient-elle, et comment est-elle
gouvernée~?
\end{interviewq}
